%! Piecewise Linear Learning Functions — Unit 8, Lesson 8.1 — Homework
\documentclass[10pt]{article}
\usepackage{linalg-article}
\usepackage{linalg-boxes}
\newcommand{\TallMath}[1]{$\displaystyle #1\rule[-1.4em]{0pt}{3.2em}$}
\newcommand{\vv}{\mathbf{v}}
\newcommand{\bb}{\mathbf{b}}
\newcommand{\zero}{\mathbf{0}}
\newcommand{\relu}{\mathrm{ReLU}}

\begin{document}

\pageheader{Unit 8, Lesson 8.1}{Homework}
\namedateperiod

\vspace{0.12in}

\begin{notesbox}{Practice}
A \textbf{piecewise-linear} function is a sum of shifted ReLUs. Each $\relu(x-s)$ adds a \textbf{fold} at
$x=s$; its coefficient sets how much the slope changes there. On any piece the slope is the sum of the
weights of the ReLUs already switched on, and $N$ folds give $N+1$ straight pieces.

\begin{enumerate}[leftmargin=1.9em, itemsep=11pt]
  \item \textbf{Shifted ReLUs.} Evaluate each, and give its fold location.
    \begin{enumerate}[label=(\alph*), leftmargin=1.8em, itemsep=4pt]
      \item $\relu(x-2)$ at $x=0$, $x=2$, $x=4$ \hfill fold at $x=\underline{\hspace{1.2cm}}$
      \item $\relu(x-3)$ at $x=1$, $x=5$ \hfill fold at $x=\underline{\hspace{1.2cm}}$
    \end{enumerate}
    \vspace{0.7cm}
  \item \textbf{Read a piecewise-linear function.} Let $F(x)=\relu(x)-\relu(x-2)+\relu(x-4)$. Compute
        $F(0)$, $F(2)$, $F(4)$, and $F(6)$. Where are the folds, what is the slope on each piece, and
        describe the shape.
        \writelines{3}
  \item \textbf{Build a tent.} Let $F(x)=\relu(x)-2\,\relu(x-3)+\relu(x-6)$. Compute $F(3)$, $F(6)$, and
        $F(9)$. Name the peak and describe the shape (what happens after $x=6$?).
        \writelines{2}
  \item \textbf{Fit data, then score it.} Use $F(x)=\relu(x)-\relu(x-2)$ as a model. Compute the
        predictions $F(1)$, $F(3)$, $F(5)$. The true targets are $1,2,3$. Compute the \textbf{loss} (sum of
        squared errors). Does $F$ fit the first two points exactly? \writelines{2}
  \item \textbf{Explain.} (a) Inside a layer, how does the \emph{bias} control where a fold sits?
        (b) Why does adding more ReLUs (a \emph{wider} weight matrix) let a network fit more complicated
        data? \writelines{2}
\end{enumerate}
\end{notesbox}

\begin{extensionbox}
\textbf{How many folds does a wiggle need?}
\begin{enumerate}[label=(\alph*), leftmargin=1.8em, itemsep=5pt]
  \item A hidden layer has $6$ ReLUs. What is the greatest number of straight pieces its function can have?
  \item A data pattern changes direction $3$ times (up, down, up, down). What is the fewest number of folds
        a piecewise-linear function needs to trace it? Explain the ``one fold per turn'' idea in a sentence.
\end{enumerate}
\writelines{3}
\end{extensionbox}

\begin{spiralbox}
\textbf{Coming next --- Lesson 8.2: Convolutional Neural Nets.} We keep the layer $\relu(A\vv+\bb)$, but for
images we \emph{reuse} the same small set of weights across the whole picture --- a special weight matrix
that scans for the same pattern everywhere.
\end{spiralbox}

\end{document}
